% introduccion.tex
Este informe presenta un análisis experimental enfocado en dos problemas fundamentales de las ciencias de la computación: el ordenamiento de arreglos unidimensionales y la multiplicación de matrices cuadradas. El objetivo principal de este estudio es relacionar el comportamiento práctico de diversos algoritmos con su respectiva complejidad teórica asintótica. Mediante la evaluación rigurosa de los tiempos de ejecución y la huella de memoria bajo distintos escenarios, tamaños y topologías de datos, se busca evidenciar cómo factores del mundo real —tales como las constantes operativas ocultas y la gestión de recursos del sistema— impactan el rendimiento final de las implementaciones teóricamente óptimas.